\section*{Chapter \ref{ch:b2:ism-star-formation} --- The Interstellar Medium and Star Formation}

\begin{solution}{exo:b2:ism-star-formation:1}
$n = \num{2.7e-24}/\num{1.67e-24} = 1.6$ hydrogen atoms per cubic centimetre. In air,
$n = P/k_{\mathrm B}T = 101325/(\num{1.381e-23}\times288) = \qty{2.55e25}{m^{-3}}$,
that is \num{2.5e19} molecules per cubic centimetre: about $1.6\times10^{19}$ times
more.
\end{solution}

\begin{solution}{exo:b2:ism-star-formation:2}
$M_J \propto n^{-1/2}$: dividing $n$ by 100 multiplies $M_J$ by 10, so
$M_J = 54.6$ solar masses at $10^2$ particles per cubic centimetre. The cloud holds
$\num{24000}/54.6 \approx 440$ Jeans masses: it cannot be stable as a whole, and must
fragment.
\end{solution}

\begin{solution}{exo:b2:ism-star-formation:3}
$c_s = \sqrt{\num{1.381e-23}\times10/(2.33\times\num{1.661e-27})} =
\qty{189}{m.s^{-1}}$ at \qty{10}{K}; $\sqrt{10}$ times more, \qty{597}{m.s^{-1}},
at \qty{100}{K}. Crossing \qty{0.1}{pc}: $\num{3.086e15}/189 = \qty{1.63e13}{s}$,
that is \qty{0.52}{Myr}.
\end{solution}

\begin{solution}{exo:b2:ism-star-formation:4}
$\rho = 2.33\times\num{1.661e-27}\times n$ with $n$ in \unit{m^{-3}}. For $10^5$
particles per cubic centimetre, $\rho = \qty{3.87e-16}{kg.m^{-3}}$ and
$t_{\mathrm{ff}} = \sqrt{3\pi/(32G\rho)} = \qty{3.38e12}{s} = \qty{0.107}{Myr}$; for
$10^2$, $\rho$ is \num{1000} times smaller and $t_{\mathrm{ff}}$ $\sqrt{1000} = 31.6$
times longer, \qty{3.38}{Myr}. The \omterm{def:b2:ism-star-formation:molecular-cloud}{dense core} collapses 32 times faster.
\end{solution}

\begin{solution}{exo:b2:ism-star-formation:5}
With $k_{\mathrm B}T/(\mu m_u) = c_s^2$, $(5k_{\mathrm B}T/G\mu m_u)^{3/2} =
5^{3/2}c_s^3G^{-3/2}$, which gives the formula. Since $M_J \propto \rho^{-1/2}$,
halving $M_J$ needs a density four times higher. Since $M_J \propto T^{3/2}$,
warming from \qty{10}{K} to \qty{20}{K} multiplies it by $2^{3/2} = 2.83$.
\end{solution}

\begin{solution}{exo:b2:ism-star-formation:6}
$t_{\mathrm{KH}} = GM^2/(RL) = \num{6.674e-11}\times(\num{1.989e30})^2/(2\times
\num{6.96e8}\times2\times\num{3.83e26}) = \qty{2.48e14}{s}$, that is
\qty{7.8}{Myr}. It is longer than the lifetime of most discs (half are gone after
about \qty{3}{Myr}): the star is still contracting when its disc disappears.
\end{solution}

\begin{solution}{exo:b2:ism-star-formation:7}
At \qty{10}{K}, $c_s^3/G = 189^3/\num{6.674e-11} = \qty{1.01e17}{kg.s^{-1}}$, that is
$\num{1.01e17}\times\num{3.156e7}/\num{1.989e30} = \num{1.6e-6}$ solar masses per
year, and one solar mass takes \qty{0.62}{Myr}. At \qty{20}{K} the rate is
$2^{3/2} = 2.83$ times higher, \num{4.5e-6} solar masses per year, and one solar
mass takes \qty{0.22}{Myr}.
\end{solution}

\begin{solution}{exo:b2:ism-star-formation:8}
$r = \sqrt{L/(16\pi\sigma)}/T^2$ with $\sqrt{\num{3.83e26}/(16\pi\times\num{5.67e-8})}
= \qty{1.159e16}{m.K^2}$: at \qty{1300}{K}, $r = \qty{6.86e9}{m} = \qty{0.046}{au}$,
well inside the orbit of Mercury. At \qty{5.2}{au}, $T = 278/\sqrt{5.2} = \qty{122}{K}$.
\end{solution}

\begin{solution}{exo:b2:ism-star-formation:9}
$v = \sqrt{GM/r} = \sqrt{\num{6.674e-11}\times1.3\times\num{1.989e30}/\num{1.496e13}}
= \qty{3.40}{km.s^{-1}}$, and $P = 2\pi r/v = \qty{2.77e10}{s} = 877$ years
(equivalently $\sqrt{100^3/1.3}$ years). In a million years the gas has made
$10^6/877 \approx \num{1100}$ turns: the disc is old compared with its own rotation,
which is why its rings are smooth circles.
\end{solution}

\begin{solution}{exo:b2:ism-star-formation:10}
$M = \int_0^R 4\pi r^2\rho_0R/r\,\mathrm dr = 2\pi\rho_0R^3$, and inside $r$,
$m(r) = 2\pi\rho_0Rr^2$. With $\mathrm dm = 4\pi\rho_0Rr\,\mathrm dr$,
\[
  U = -\int_0^R \frac{G\,m(r)}{r}\,\mathrm dm
    = -8\pi^2G\rho_0^2R^2\int_0^R r^2\,\mathrm dr = -\frac{8\pi^2}{3}G\rho_0^2R^5
    = -\frac23\,\frac{GM^2}{R},
\]
using $\rho_0 = M/(2\pi R^3)$. It is $\tfrac{2/3}{3/5} = 1.11$ times the uniform value:
the concentrated cloud is more tightly bound and needs \qty{11}{\percent} more
thermal energy to be held up; it collapses more easily.
\end{solution}

\begin{solution}{exo:b2:ism-star-formation:11}
Period $2\pi/\Omega = \qty{6.28e14}{s}$, about \qty{20}{Myr}. With
$R = \qty{3.086e15}{m}$, $j = \Omega R^2 = \qty{9.52e16}{m^2.s^{-1}}$, and
$r_c = j^2/(GM) = \num{9.07e33}/\num{1.327e20} = \qty{6.83e13}{m} = \qty{457}{au}$.
In solid rotation $j = \Omega r^2$ grows as the square of the distance from the axis:
gas from inside the core has less angular momentum, and lands closer to the star.
\end{solution}

\begin{solution}{exo:b2:ism-star-formation:12}
$\bar t = 2.46$, $\bar f = 0.588$; $\sum(t - \bar t)^2 = 12.612$,
$\sum(t - \bar t)(f - \bar f) = -2.018$, so $b = 2.018/12.612 = 0.160$ per million
years and $a = 0.588 + 0.160\times2.46 = 0.982$. The line reaches $0.5$ at
$(0.982 - 0.5)/0.160 = \qty{3.0}{Myr}$ and zero at $0.982/0.160 = \qty{6.1}{Myr}$, as in
the figure. A fraction cannot fall below zero: older clusters, such as NGC 1960 at
\qty{30}{Myr}, keep a small fraction near zero that a straight line cannot describe.
\end{solution}

\begin{solution}{pb:b2:ism-star-formation:1}
\textbf{1.} $R = 0.08\times\num{3.086e16} = \qty{2.469e15}{m}$.

\textbf{2.} $N = \qty{1.5e26}{m^{-2}}$, so $2.8\times\num{1.661e-27}\times\num{1.5e26} =
\qty{0.698}{kg.m^{-2}}$.

\textbf{3.} $M = 0.698\times\pi R^2 = 0.698\times\num{1.915e31} = \qty{1.34e31}{kg}$,
that is $6.7$ solar masses.

\textbf{4.} $V = \tfrac43\pi R^3 = \qty{6.30e46}{m^3}$ and
$\rho = M/V = \qty{2.12e-16}{kg.m^{-3}}$.

\textbf{5.} $n = \rho/(2.33m_u) = \qty{5.48e10}{m^{-3}}$, that is \num{5.5e4} particles
per cubic centimetre.

\textbf{6.} Yes: its density lies between $10^4$ and $10^5$, its radius between 0.03 and
\qty{0.2}{pc}; its mass, 6.7 solar masses, is at the top of the range of cores.

\textbf{7.} $M_J = 5.46\times(10^4/\num{5.48e4})^{1/2} = 2.33$ solar masses.

\textbf{8.} $6.7/2.33 = 2.9$ Jeans masses: the core is unstable and collapses.

\textbf{9.} $R_J = 0.1316\times(10^4/\num{5.48e4})^{1/2} = \qty{0.056}{pc}$, smaller than
the radius, \qty{0.08}{pc}: the same conclusion.

\textbf{10.} $2^{3/2}\times2.33 = 6.6$ solar masses: about the core's mass. At
\qty{20}{K} the core would be marginal; cold is what makes it collapse.

\textbf{11.} $t_{\mathrm{ff}} = \sqrt{3\pi/(32\times\num{6.674e-11}\times\num{2.12e-16})}
= \qty{4.56e12}{s}$, about \num{145000} years.

\textbf{12.} $R/c_s = \num{2.469e15}/189 = \qty{1.31e13}{s}$, \num{414000} years: three
times longer than the free fall. A pressure wave cannot cross the core before it has
collapsed; pressure cannot react fast enough to stop the collapse.

\textbf{13.} $c_s = \qty{189}{m.s^{-1}}$ (\cref{exo:b2:ism-star-formation:3}).

\textbf{14.} $c_s^3/G = \qty{1.01e17}{kg.s^{-1}}$, \num{1.6e-6} solar masses per year.

\textbf{15.} $1/\num{1.6e-6} = \num{624000}$ years.

\textbf{16.} A third of 6.7 is $2.2$: two stars of one solar mass.

\textbf{17.} $t_{\mathrm{KH}} = GM_\odot^2/(3R_\odot\times4L_\odot) = \qty{8.25e13}{s}$,
\qty{2.6}{Myr}.

\textbf{18.} Longer, four times the accretion time: when accretion stops, the star is
still contracting, shining without fusion; it is a \omterm{def:b2:ism-star-formation:t-tauri}{T Tauri star}, not yet on the main
sequence.

\textbf{19.} $\Omega = \qty{0.4e3}{m.s^{-1}}/\qty{3.086e16}{m} = \qty{1.30e-14}{rad.s^{-1}}$.

\textbf{20.} $j = 1.30\times10^{-14}\times(\num{2.469e15})^2 = \qty{7.90e16}{m^2.s^{-1}}$.

\textbf{21.} $r_c = j^2/(GM_\odot) = \num{6.24e33}/\num{1.327e20} = \qty{4.70e13}{m}$, that
is \qty{314}{au}: larger than the HL Tauri disc, as expected for gas from the edge of the
core.

\textbf{22.} $T = (L/16\pi\sigma r^2)^{1/4} = \qty{278}{K}$ at \qty{1}{au}, and
$278/\sqrt{30} = \qty{51}{K}$ at \qty{30}{au}.

\textbf{23.} $r = \num{1.159e16}/180^2 = \qty{3.58e11}{m}$, that is \qty{2.4}{au}, as in the text.

\textbf{24.} The star of one solar mass takes
\textbf{\qty[text-series-to-math=true]{624}{} thousand years} to assemble, about four
times the \omterm{def:b2:ism-star-formation:free-fall}{free-fall time} of its core (\num{145000} years): the outer gas reaches the
\omterm{def:b2:ism-star-formation:protostar}{protostar} only gradually, inside out.
\end{solution}
